\section{讨论}

\subsection{方法论贡献}

本文提出了一种基于双时间尺度的分层强化学习方法，用于解决动态热点场景下多无人机协同计算卸载问题。该方法的核心贡献体现在以下三个方面。

\textbf{双时间尺度决策架构。} 传统方法通常在单一时间尺度上处理所有决策问题，或者简单划分模块而忽略宏观与微观决策的时域差异。本文首次将PPO宏观轨迹控制与SAC微观路由决策有机结合：宏观层以大周期为粒度决定无人机质心部署，低层以物理步为粒度执行具体的任务路由选择。这种分离设计使两层可以针对各自时间尺度的动态特性进行优化，而非被迫在冲突的时间粒度上妥协。

\textbf{异频更新机制。} 宏观PPO控制器以300个宏步为周期进行策略更新，保持策略的稳定性；微观SAC以20帧为单位进行快速更新，能够快速响应通信条件的变化。这种异步更新频率的设计借鉴了双时间尺度优化的理论直觉：慢速更新的上层提供稳定的行为指导，快速更新的下层在给定指导下进行细粒度优化。实验表明，这种更新频率的差异对于训练稳定性至关重要，过快或过慢的上层更新都会导致性能下降。

\textbf{分阶段训练策略。} 两层强化学习的联合训练一直是HRL领域的核心难题，原因在于两层策略相互依赖、梯度耦合，容易导致训练崩溃或振荡。本文采用质心预测Oracle引导的预训练策略，先让下层SAC学习在理想部署下的最优路由，再逐步引入PPO学习部署策略。这种课程学习式的分阶段训练有效降低了训练复杂度，使原本难以稳定的联合优化问题变得可解。

\subsection{与现有工作的区别}

将本文方法与相关工作进行对比，可以更清晰地理解其创新定位。

\textbf{与Xu等~\citep{xu2022trajectory}的比较。} Xu等提出的方法在微观层采用凸优化进行静态路由规划，将计算卸载问题转化为可精确求解的数学规划问题。这种方法在通信拓扑相对稳定的场景下表现良好，但其局限性在于：无法适应通信条件的快速变化，难以处理随机到达的动态任务，且无法从与环境的交互中持续改进。本文采用学习型微观协作策略，用SAC替代静态优化，使系统能够在运行时学习并适应复杂的无线信道特性和任务到达模式。更重要的是，学习型方法可以泛化到凸优化难以处理的非凸场景或约束复杂的情况。

\textbf{与传统HRL方法的比较。} 传统HRL方法如 feudal reinforcement learning~\citep{dayan1993feudal}、HIRO~\citep{nachum2018data}等虽然也采用层次结构，但通常隐式地分离不同粒度的决策，未明确利用宏观/微观的物理时间语义。本文显式利用了任务路由与轨迹控制在物理时间尺度上的自然分离，使宏观决策以"驻留大周期"为单位，微观决策以"物理步"为单位。这种设计不仅在理论上更符合问题的物理本质，在实践中也取得了更好的性能。

\textbf{与联合优化方法的比较。} 端到端联合优化一直是强化学习的理想目标，但其训练难度往往被低估。将部署决策和路由决策视为同一马尔可夫决策过程的两个动作维度会导致状态-动作空间急剧膨胀，且两层策略的梯度相互干扰，严重影响训练稳定性。本文通过模块化解耦和分阶段训练策略，在保持一定协同优化的同时显著降低了训练难度。实验表明，这种解耦并未牺牲太多性能，反而因为各模块能够专注于各自的子问题而获得了更好的最终表现。

\subsection{理论分析}

本文方法的理论有效性可以从以下几个角度理解。

\textbf{双时间尺度近似的合理性。} 部署决策（如无人机质心位置）相对于任务到达和路由决策具有较慢的动态特性。这种时间尺度分离在大规模通信系统中是常见的~\citep{neely2005stochastic}。本文方法实质上是对这种分离的显式利用，而非人为强加。当下层SAC在给定部署下优化路由时，上层PPO可以在更长的时间窗口内评估部署策略的效果，从而做出更具前瞻性的决策。

\textbf{统一折扣框架的意义。} 传统的层次强化学习往往为不同层次设计独立的奖励函数或使用不同的折扣因子，这会导致各层优化目标不一致。本文采用统一的物理折扣因子处理两层反馈，使宏观部署的长期效果能够被正确传递到微观决策。这种统一处理虽然增加了训练的挑战性，但保证了整个系统以一致的目标为导向。

\textbf{执行路由表征的一致性。} 本文在执行路由评论家中保持与策略动作编码的一致性，即评论家输出的是与执行动作相对应的价值估计。这一设计选择虽然限制了评论家的表达能力，但保证了价值估计的语义正确性。实验中的消融研究验证了这一设计的重要性：使用宽松动作空间的评论家虽然能获得更高的价值估计，但策略性能反而下降。

\subsection{局限性与未来工作}

尽管本文方法在仿真环境中取得了显著的性能提升，但仍存在以下局限性，值得在未来研究中深入探索。

\textbf{场景泛化能力。} 本文的所有实验均在开发场景（源区域数量、用户密度、无人机数量等）上进行验证。在实际部署中，场景参数可能与训练环境存在显著差异。例如，无人机数量可能动态变化，用户分布可能呈现不同的空间相关性，无人机硬件能力可能存在异构性。未来工作需要在更多样的场景下验证方法的可迁移性，特别是从密集场景到稀疏场景、从同构无人机到异构无人机的迁移能力。

\textbf{Oracle追踪能力的提升。} PPO控制器追踪Predictive Oracle的能力仍然有限。虽然600窗口扩展训练后差距有所缩小，但仍未完全收敛。这可能是因为PPO的探索策略与Oracle的确定性规划之间存在本质差异。未来的改进方向包括：引入计划预测模块作为额外输入、采用基于模型的规划方法辅助策略学习、或设计专门的课程学习方案使PPO逐步逼近Oracle的性能。

\textbf{联合训练的增量价值。} 当前实验尚未充分验证联合训练的增量价值。独立训练的质心部署加SAC基线与完整方法之间的差距需要更严格的消融实验来量化。未来应设计对照实验，在相同的计算预算下比较联合训练与独立训练的效果差异，以明确两层协同学习的实际贡献。

\textbf{实际部署的工程考量。} 本文主要关注算法层面的设计与验证，在实际部署中还需考虑以下工程因素：通信延迟可能导致状态观测的异步性；无人机电池容量限制了悬停时间和计算资源；机间通信的可靠性受天气、障碍物等因素影响；无人机硬件故障或任务优先级变化需要在线适应机制。将仿真验证拓展到半实物仿真或实际飞行测试是未来工作的重要方向。

\textbf{多无人机协同的规模化。} 当无人机数量进一步增加时，现有架构可能面临协调复杂度的挑战。大规模多智能体强化学习中的通信、credit分配、训练稳定性等问题需要专门的技术方案。此外，无人机之间的竞争与合作关系也需要更精细的建模。

\textbf{安全与鲁棒性保证。} 强化学习策略在部署时可能面临分布外（out-of-distribution）情况，其行为难以预测。本文未涉及安全约束的显式建模，如禁飞区避让、碰撞避免、能耗下限保障等。将安全强化学习或约束强化学习方法融入当前架构是保障实际安全部署的关键。

综上所述，本文提出的双时间尺度分层强化学习方法为动态热点场景下的多无人机协同计算卸载问题提供了有竞争力的解决方案，但其全部潜力仍有待在更广泛的场景和更严格的验证中得到充分释放。
